\begin{theorem}[Rado embedding]
Let $r$ be a quasi-order on $Q$.  If $r$ is well-quasi-ordered and the
first-level domination relation $\mathsf{SetDomination}(r)$ on
$\mathcal P(Q)$ is not well-quasi-ordered, then there is a map

\[
e:\mathsf{IncreasingPair}\longrightarrow Q
\]

such that $e$ is a relation embedding from the Rado order to $r$:

\[
\mathsf{RelationEmbedding}(\mathsf{RadoOrder},r,e).
\]
Equivalently, $e$ is injective and, for all increasing pairs $p,q$,
\[
\mathsf{RadoOrder}(p,q)\quad\Longleftrightarrow\quad r(e(p),e(q)).
\]
\end{theorem}
\begin{proof}
By the bad-witness characterization \noderef{PowerWqoBadPairSequence},
the hypothesis that $\mathsf{SetDomination}(r)$ is not well-quasi-ordered
gives a map $f:\mathsf{IncreasingPair}\to Q$ satisfying
\noderef{BadPairSequence}:
\[
 m<n<l\quad\Longrightarrow\quad
 \neg r\bigl(f(m,n),f(n,l)\bigr).
\]
We use the two finite homogeneous consequences of Nash--Williams,
namely \noderef{TripleHomogeneous} and
\noderef{QuadrupleHomogeneous}, exactly as in the paper.

Color every three-element subset $\{i,j,k\}$, written with
$i<j<k$, by $1$ when
\[
r\bigl(f(i,j),f(i,k)\bigr)
\]
holds, and by $0$ otherwise.  Apply \noderef{TripleHomogeneous} to the
whole set of natural numbers.  We obtain an infinite set $N$ on which
the color is constant.  The constant color cannot be $0$: fix any
$i\in N$ and enumerate the infinite tail $N/i$ increasingly, as allowed
by \noderef{InfiniteSetEnumeration}.  The resulting sequence
$t\mapsto f(i,n_t)$ has no forward related pair, because its two terms
at indices $u<v$ occur in the colored triple
$i<n_u<n_v$.  This contradicts $h_Q$, since
$\mathsf{WellQuasiOrdered}(r)$ means that every sequence has indices
$u<v$ with related values.  Thus every triple in $N$ has color $1$:
\[
i<j<k\text{ in }N\quad\Longrightarrow\quad
r\bigl(f(i,j),f(i,k)\bigr).
\tag{1}
\]

Now color every four-element subset $\{i,j,k,l\}$ of $N$, with
$i<j<k<l$, by $1$ when
\[
r\bigl(f(i,j),f(k,l)\bigr)
\]
holds, and by $0$ otherwise.  Applying
\noderef{QuadrupleHomogeneous} gives an infinite $M\subseteq N$ on
which this color is constant.  The constant cannot be $0$.  Indeed,
use \noderef{InfiniteSetEnumeration} to enumerate $M$ increasingly as
$z$.  Taking $i_u=z(2u)$ and $j_u=z(2u+1)$ gives
\[
i_0<j_0<i_1<j_1<i_2<j_2<\cdots
\]
because $z(2u)<z(2u+1)<z(2u+2)$.  If the color were $0$, then for
every $u<v$ the quadruple
$i_u<j_u<i_v<j_v$ would give
\[
\neg r\bigl(f(i_u,j_u),f(i_v,j_v)\bigr),
\]
so the sequence $u\mapsto f(i_u,j_u)$ would contradict $h_Q$.  Hence
\[
i<j<k<l\text{ in }M\quad\Longrightarrow\quad
r\bigl(f(i,j),f(k,l)\bigr).
\tag{2}
\]

Let $a=\min M$ and let $X=M\setminus\{a\}$.  The set $X$ is
infinite: otherwise $M=X\cup\{a\}$ would be finite, contrary to the
construction of the infinite set $M$.  Use
\noderef{InfiniteSetEnumeration} to choose a strictly increasing
enumeration $x:\mathbb N\to\mathbb N$ whose range is exactly $X$.
Thus $x(i)\in X$ for every $i$, and $x(i)<x(j)$ whenever $i<j$.

Before defining the embedding, we record the direct monotonicity fact
that will be used throughout.  If $u<v$ and $s<t$ are natural numbers
in $M$ and
\[
\mathsf{RadoOrder}((u,v),(s,t)),
\]
then
\[
r(f(u,v),f(s,t)).
\tag{3}
\]
Indeed, by the two alternatives in \noderef{RadoOrder}, either
$u=s$ and $v\leq t$, or $v<s$.  In the first alternative, if $v=t$,
reflexivity of the quasi-order gives the conclusion.  If $v<t$, then
$u<v<t$ are in $M$, so the triple homogeneity relation (1) gives
\[
r(f(u,v),f(u,t)),
\]
which is the required conclusion because $s=u$.
In the second alternative, $u<v<s<t$ are in $M$, so the quadruple
homogeneity relation (2) gives the conclusion directly.  This proves
(3), including both clauses of the Rado order, before any restriction
to $X$ is made.

For an increasing pair $p=(i,j)$, define
\[
e(p)=f\bigl(x(i),x(j)\bigr).
\]
The strict increase of $x$ makes $x(i)<x(j)$, so this is a valid
argument of $f$, and \noderef{BadPairSequence} applies to all such
values.

We first prove preservation of the Rado order on these pairs.  Let
$p=(i,j)$ and $q=(k,l)$.  If the first clause of
\noderef{RadoOrder} holds, then $i=k$ and $j\leq l$, hence
$x(i)=x(k)$ and $x(j)\leq x(l)$; the direct fact (3) gives
$r(e(p),e(q))$.  If the second clause holds, then $j<k$, and strict
increase gives
$x(i)<x(j)<x(k)<x(l)$; again (3) gives $r(e(p),e(q))$.  Therefore
\[
\mathsf{RadoOrder}(p,q)\Longrightarrow r(e(p),e(q)).
\tag{4}
\]

For reflection, assume $r(e(p),e(q))$, where $p=(i,j)$ and
$q=(k,l)$.  Suppose that $\mathsf{RadoOrder}(p,q)$ fails.  By the two
clauses in \noderef{RadoOrder}, this means
\[
k\leq j\quad\text{and}\quad (i\neq k\ \text{or}\ l<j).
\tag{5}
\]
There are three cases.  If $l<j$, take $n=j+1$, so $j<n$ and
$x(j)<x(n)$.  The second clause of \noderef{RadoOrder} gives
\[
\mathsf{RadoOrder}((x(k),x(l)),(x(j),x(n))).
\]
Applying the direct M-monotonicity fact (3) to these literal pairs gives
\[
r\bigl(f(x(k),x(l)),f(x(j),x(n))\bigr).
\]
Composing this with the assumed relation
$r(f(x(i),x(j)),f(x(k),x(l)))$ and using transitivity gives
\[
r\bigl(f(x(i),x(j)),f(x(j),x(n))\bigr),
\]
contradicting \noderef{BadPairSequence} at the increasing triple
$x(i)<x(j)<x(n)$.

Assume next that $l\geq j$ and $i<k$.  From (5), $k\leq j$, so the
first clause of \noderef{RadoOrder}, together with strict increase of
$x$, gives
\[
\mathsf{RadoOrder}((x(i),x(k)),(x(i),x(j))).
\]
Applying (3) to these literal pairs gives
\[
r\bigl(f(x(i),x(k)),f(x(i),x(j))\bigr).
\]
Composing with the assumed relation gives
\[
r\bigl(f(x(i),x(k)),f(x(k),x(l))\bigr),
\]
which contradicts \noderef{BadPairSequence} at
$x(i)<x(k)<x(l)$.

The remaining possibility in (5) is $l\geq j$ and $k<i$.  The minimum
$a$ of $M$ is strictly below every element of $X$, so in particular
$a<x(k)$ even when $k=0$.  Since $k<i$, strict increase gives
$x(k)<x(i)$, and therefore the literal increasing pairs
$(a,x(k))$ and $(x(i),x(j))$ satisfy the second clause of
\noderef{RadoOrder}.  The direct fact (3) gives
\[
r\bigl(f(a,x(k)),f(x(i),x(j))\bigr).
\]
Composing with the assumed relation gives
\[
r\bigl(f(a,x(k)),f(x(k),x(l))\bigr),
\]
contradicting \noderef{BadPairSequence} at the increasing triple
$a<x(k)<x(l)$.  All cases are impossible, so
\[
r(e(p),e(q))\Longrightarrow\mathsf{RadoOrder}(p,q).
\tag{6}
\]

It remains to check injectivity.  If $e(p)=e(q)$, reflexivity of $r$
gives both $r(e(p),e(q))$ and $r(e(q),e(p))$.  By (6), both Rado
relations hold.  The two-clause definition in \noderef{RadoOrder} is
antisymmetric: a cross-clause relation $p\leq q$ cannot be reversed,
and if the first clause holds in both directions then the first
coordinates agree and antisymmetry of $\leq$ on $\mathbb N$ gives equality
of the second coordinates.  Hence $p=q$.  Thus $e$ is injective.

Combining injectivity with (4) and (6), and unfolding
\noderef{RelationEmbedding}, proves
\[
\mathsf{RelationEmbedding}(\mathsf{RadoOrder},r,e).
\]
This is the embedding asserted in lines 985--987 of the paper.
\end{proof}
